\section{Additional Implementation Details}\label{appx:distrib_var}
\label{appx:more_implementational_details}
\subsection{Pre-whitening of Data}
\label{appx:prewhiten}
For all datasets covered in \autoref{real_world_data} the following ``whitening'' preprocessing step is adopted:
Let $\vec{y}$ be the vector of all regressor values in the {\it training} dataset only, and $X$ the matrix of all regressands in the {\it training} dataset only, where each row of $X$ is a feature.
Let $\mu_y, \sigma_y^2$ be the sample mean and variance of $\vec{y}$ respectively in the training dataset, then the whitened $y$ values used in both the training and test set are simply $\sigma_y^{-1}(y - \mu_y)$.
Let $\vec{\mu}_X, \Sigma_X$ be the sample mean and covariance matrix of $X$ respectively .
Let $\Sigma_X = M M^T$, then the whitened $x$ values in both the training and test data are $\frac{1}{\sqrt{d}}M^{-1}(\vec{x} - \vec{\mu}_X)$, where $d$ is the feature dimension of $X$. \textbf{Note}: the performance metrics given in \autoref{real_world_data} are expressed in terms of the whitened $y$ values rather than the $y$ values in their original form. This appears to be common practice in the literature and has no bearing on the {\it comparative} performance of the different methods within this paper.

\subsection{Test-Set Batching}
To prevent excessive memory consumption, we perform all predictions for the distributed and variational methods in batches of 1000 points at a time. Where this is not possible (e.g. for especially large datasets), we use smaller batches of 500 or 250 points, as appropriate.

\subsection{Additional Implementation Details for SVGP}
We use the sparse variational inducing point approach of \cite{Hensman}, following the implementation provided by GPyTorch, which in particular uses a Choleksy decomposition to parameterise the covariance matrix of the variational prior.
We broadly follow the SVGP implementation example provided by \url{https://docs.gpytorch.ai/en/stable/examples/04_Variational_and_Approximate_GPs/SVGP_Regression_CUDA.html}. In particular, we follow their example in using the Adam optimiser to train our model over 100 epochs with a minibatch size of 1024 and a learning rate of 0.01. We opt to use 1024 inducing points. All experiments under this method are run on a SageMaker ml.p3.2xlarge instance, consisting of a single Tesla V100 GPU with 16GB of memory.


\subsection{Additional Implementation Details for Distributed methods}
 A good introduction to distributed methods for Gaussian process inference is \cite{deisenroth2015distributed}. Here we run the product-of-experts (PoE) \cite{hinton2002training}, generalised product-of-experts (gPoE) \cite{cao2014generalized}, Bayesian committee machine (BCM) \cite{tresp2000bayesian}, robust Bayesian committee machine (rBCM) \cite{deisenroth2015distributed} and generalised robust Bayesian committee machine (GrBCM) \cite{amalg_SOD} following the recommendation in \cite{Cohen20b} to aggregate in $f$-space.
 There are three components to any distributed method: the hyperparameter inference, the \emph{partitioner} and the \emph{aggregator}.
Hyperparameter estimation is the same for all of the methods: we use the method in section 3.1 of \cite{deisenroth2015distributed}, randomly partitioning the entire training set into subsets of size 625 (or as close as possible with equal-sized experts given that in general $n$ is not a multiple of 625). A block diagonal approximation (with $n/625$ blocks) is then used to approximate to the full $n \times n$ gram kernel matrix. To recover hyperparameters with this we use Gaussian Process models with a zero prior mean and a scaled square-exponential kernel. Training is conducted using the Adam optimiser with a learning rate of 0.1 over 100 optimiser iterations. Once the hyperparameters are trained, we run our distributed prediction mechanism to evaluate performance against the test-set. The 625-sized partitioned blocks are referred to as ``experts'' and the shared hyperparameter values are distributed to each expert and held fixed thereafter. In the \emph{aggregator}, or distributed prediction phase, each expert produces an individual predictive distribution and these are then aggregated to a final predictive mean and variance for each of our test points. GRBCM prediction is a little more complex than this as it makes use of an additional ``communications'' expert as explained in \cite{amalg_SOD}, aggregating in $f$-space as recommended in \cite{Cohen20b}. We provide timing statistics for training these models.

We use our own GPyTorch-based implementation of distributed GP approximations.
All exact GP calculations are performed using GPyTorch using the default settings (so 20 Lanczos iterations throughout and a CG tolerance of 1 for hyperparameter inference, and $10^{-3}$ for posterior predictions).
For all of our experiments, we utilise an AWS t3.2xlarge instance (consisting of 8 Intel Skylake Processors and 32 GB of RAM).

\subsection{Reproducibility}
All code used to generate tables and figures in this document can be found at \url{https://github.com/ant-stephenson/gpnn-experiments/}.



